% SPDX-License-Identifier: Apache-2.0
% covers: Tee
\datasheet{Tee}{One stream to two receivers in step: a word is taken when both have room, and neither can run a word ahead of the other.}

\dsfacts{\code{lib/parts/src/redundant.rs}}{\code{txhdl\_parts::redundant::Tee}}%
{\code{//docs:vreteno}}

\subsection*{Function}

A machine that must not be quietly wrong runs two of itself and
compares. The comparison is worth something only if both copies see
exactly the same words in the same cycles, and \code{Tee<T>} is what
makes that true of a channel: it takes one word from its input and
offers it to both sides, and it takes the next only once both have had
it.

That is the whole of the part. It does not compare, drop or reorder
anything; \code{Check} is the other half.

\subsection*{Parameters}

\code{T}, the payload, which must be a \code{Transaction} and a
\code{Value}. The tables use the tree's read answer,
\code{R<32, 2>}. \code{Pair} also tees \code{Done<2>} and
\code{Grant<2>}, and the example and the tests tee a byte.

\dstables{Tee}

\subsection*{Behaviour}

\begin{itemize}
\item The word is registered rather than passed through. A unit that
read a channel and drove one in the same cycle would be a
combinational path from one side's handshake to the other's, which is
the shape this repository keeps out of its channels, and its netlist
could not be replayed against the run cycle by cycle. So a word taken
in one cycle is offered from the next.
\item \code{full} says a word is held. Each side takes it in the first
cycle in which that side has room and has not taken it already, which
\code{took\_a} and \code{took\_b} record.
\item A new word is taken when the held one has been taken by both, or
when none is held. So the tee refuses its input for as long as the
slower side has not caught up, and the faster side cannot be offered
the next word early.
\item With both sides ready in every cycle the tee passes one word a
cycle, and the two sides are always on the same word.
\end{itemize}

\subsection*{Verification}

\begin{itemize}
\item \code{both\_sides\_of\_a\_tee\_see\_every\_word} and
\code{a\_slow\_side\_holds\_the\_other\_back}, in
\code{//lib/parts:parts\_test}. The second gives one side a receiver
that takes a word every other cycle and finds the two sides' lists
equal, which is the property the part exists for.
\item \code{//lib/examples:ex\_redundant} runs two adders behind one
tee and gives one of them a different addend halfway through. The build
simulates the netlist against that run under nvc and Verilator:
\code{//docs:redundant\_sim\_tee\_tee\_tb\_test} and
\code{//docs:redundant\_vsim\_tee\_test}.
\end{itemize}

\subsection*{Used by}

\code{Pair}, three times: the read answers, the write completions and
the identifiers the tracker hands out, each to both cores.

\subsection*{Limits}

Two receivers, not \code{N}. One word of storage, so one cycle between
taking a word and offering it, and no queue behind that. A side that
never becomes ready stops the tee, which is the intended behaviour and
also means the part has no timeout. There is no reset port: the
registers start at their defaults.
